\section{Asymptotic Properties}
\label{sec:theory_text}
In \Cref{sec:theory_simple} we focused on a low-overlap setting with $P[1/\pi(X)]=\infty$ where true propensities are known. In contrast, here we consider a more standard setting and identify conditions under which  the C-Learner is semiparametrically efficient and doubly robust. In particular, our theoretical treatment also shows the asymptotic optimality of one-step estimation (self-normalized AIPW) and targeting (TMLE with unbounded real outcomes) since they are also C-Learners satisfying the requisite conditions. 

The proof techniques here are common across those used for other first-order de-biasing methods~\citep{van2006targeted,van2011targeted,van2016one,van2017generally,Kennedy22} as we share the goal of setting the first-order error term in the distributional Taylor expansion to zero, and also showing second-order terms are negligible. 


As in \Cref{sec:theory_simple}, we focus on the cross-fitted~\citep{van2011cross,ChernozhukovChDeDuHaNeRo18} formulation of the C-Learner in \Cref{sec:data_splitting}. %
In addition to its practical benefits,
cross-fitting simplifies the proof of  asymptotic optimality, especially when nuisance parameter models can be large and complex. When nuisance  model classes are Donsker---converging at $\sqrt{n}$-rates---asymptotic optimality can be shown without explicit sample splitting~\citep{Kennedy22,vaart2023empirical,van2000asymptotic}.
The constraint~\eqref{eq:constraint} is over $P_{k,n}$. %

\ifdefined\response\newpage\newpage\clearpage\pagebreak\else\fi
\label{sec:ass_d_note}
We shortly identify sufficient conditions that guarantee asymptotic optimality  (Assumption~\ref{ass:empirical_proc_mean}), in addition to standard conditions on the nuisance parameters $\what{\pi}_{-k,n},\what{\mu}^C_{-k,n}$ (\Cref{ass:overlap}, \Cref{ass:consistency}) and on outcomes (\Cref{ass:bounded_outcomes}). 
\begin{assumption}[Overlap]
\label{ass:overlap}
    For some $\eta>0$ and for all $k,n$, we have  
    \vspace{-0.2em}
    $$
    \eta\leq  \pi(X) \le 1-\eta, \quad \eta \leq \what \pi_{-k,n}(X) \leq 1-\eta \quad \text{a.s.}
    $$  
\end{assumption}
\vspace{-1.2em}
\begin{assumption}[Convergence rates of propensity and constrained outcome models]
\label{ass:consistency}
For all~$k \in \{1,\ldots,K\}$, 
both $\what\pi, \what\mu^C$ are consistent,\vspace{-0.3em}
    $$
    \|\what\pi_{-k,n} - \pi\|_{L_2(P)}=o_P(1), \quad  \|\what\mu_{-k,n}^C - \mu\|_{L_2(P)}=o_P(1)
    $$ 
$$\text{and also} \quad    \|\what\pi_{-k,n} - \pi\|_{L_2(P)} \cdot  \|\what\mu_{-k,n}^C - \mu\|_{L_2(P)}=o_P(n^{-\frac{1}{2}}).$$
\end{assumption}
\noindent As we discuss below, we can relax these assumptions
when guaranteeing double robustness (consistency under misspecified nuisance parameters).
As is typical, we assume that outcomes do not differ too much from their means, conditional on covariates. 
\begin{assumption}
[Outcomes close to conditional means] \label{ass:bounded_outcomes} For all $k,n$ and some $0<B<\infty$, 
\vspace{-0.3em}
$$\|A(Y-P[Y\mid X])\|_{L_2(P)} =\|A(Y-\mu(X))\|_{L_2(P)} \leq B. $$
\end{assumption}

We also require the following condition (\Cref{ass:empirical_proc_mean}) on the C-Learner outcome model. This condition is new to our setting, and warrants discussion. %

\begin{assumption}[Empirical process assumption]%
\label{ass:empirical_proc_mean}\ \vspace{-0.3em}
$$(P_{k,n}-P)(  \what\mu_{-k,n}^C(X)- \mu(X)) =o_P(n^{-1/2}).$$
\end{assumption}
\noindent In \cref{sec:other_methods}, we showed how versions of one-step estimation  and targeting can also be considered C-Learners. We show in \cref{sec:show_constant_shift} and \cref{sec:show_tmle} that cross-fitted versions of the self-normalized AIPW and TMLE for continuous and unbounded outcomes, under standard assumptions on $\what\pi_{-k,n}, \what \mu_{-k,n}$ (with $\what \mu_{-k,n}$ the usual outcome model learned using Equation~\eqref{eq:regular_outcome}) also satisfy Assumption~\ref{ass:empirical_proc_mean}. Thus, C-Learner results in Theorems~\ref{thm:asymp},~\ref{thm:dr} (to come) apply directly to these estimators. %
One future research direction is to understand the model classes and constrained optimization methods for which this assumption holds. 

\Cref{ass:empirical_proc_mean} is not implied by previous assumptions. Consider the following example: 

\begin{example}[Assumption \ref{ass:empirical_proc_mean} does not follow from Assumptions \ref{ass:overlap}, \ref{ass:consistency}, \ref{ass:bounded_outcomes}]\label{ex:ass_D}
Assume \Cref{ass:overlap}.
Let $\what\mu_{-k,n}$ be an \emph{unconstrained} solution to
$\what\mu_{-k,n}\in \argmin_{\tilde\mu\in\mathcal F} P_{-k,n}[A(Y-\tilde \mu(X))^2]$,
in contrast txo the constrained problem in \eqref{eq:constraint}. 
Then, define 
$\what\mu^C_{-k,n}(x)=\what\mu_{-k,n}(x)+\sum_{i=1}^{n/K} \mathbf{1} \left\{x=x_i\right\}$, with the sum of indicators over elements in the $P_{k,n}$ fold.
Then 
\begin{align}
\label{eq:01_counterex}
    P\left|\what\mu^C_{-k,n}(X)-\what\mu_{-k,n}(X)\right|=0\;\text{ and } \;
P_{k,n}\left(\what\mu^C_{-k,n}(X)-\what\mu_{-k,n}(X)\right)=1.
\end{align}
We make the following additional assumptions for this example. 
We also assume that $\|\what\mu_{-k,n}(X)-\mu(X)\|=o_P(1)$ (analogous to \Cref{ass:consistency}).
Let $\what\mu$ denote $\argmin_{\tilde \mu\in\mathcal F} P[A(Y-\tilde\mu(X))^2]$, the best fitted model over the entire population. For simplicity, assume that $Y$ is bounded to satisfy \Cref{ass:bounded_outcomes}, and that $\what\mu(X)$ is bounded as well. Then
\vspace{-0.7em}
\begin{align*}
(P_{k,n}-P)\left(\what\mu^C_{-k,n}(X)-\mu(X)\right)&=
(P_{k,n}-P)\left(\what\mu^C_{-k,n}(X)-\what\mu_{-k,n}(X)\right)
\\&+
(P_{k,n}-P)\left(\what\mu_{-k,n}(X)-\what\mu(X)\right)\\
&+(P_{k,n}-P)\left(\what\mu(X)-\mu(X)\right)  
\end{align*}
where the first term on the RHS is 1 from \eqref{eq:01_counterex}, the second is $o_P(1)$ by the cross-fitting lemma (\Cref{lem:cross_fitting} in \Cref{sec:show_for_aipw_tmle}), and the third is $o_P(1)$ by Chebyshev. 
Therefore, 
\vspace{-0.7em}
$$(P_{k,n}-P)\left(\what\mu^C_{-k,n}(X)-\mu(X)\right)=1+o_P(1)$$
so that 
Assumption~\ref{ass:empirical_proc_mean} does not hold for this example, while Assumptions \ref{ass:overlap}, \ref{ass:consistency}, \ref{ass:bounded_outcomes} do. 
\end{example}
\label{sec:ass_d_note-end}
\ifdefined\response\newpage\newpage\clearpage\pagebreak\else\fi






\looseness=-1 The C-Learner~\eqref{eq:c_learner_def} enjoys the following asymptotics; see \cref{sec:asymptotics} for the proof.
\begin{theorem}[Asymptotic variance of C-Learner]
\label{thm:asymp}
 Under Assumptions~\ref{ass:overlap}, \ref{ass:consistency}, 
 \ref{ass:bounded_outcomes}, and
\ref{ass:empirical_proc_mean},
 \begin{equation*}
    \sqrt{n} 
    (\what{\psi}_{n}^C - \psi(P))
    \cd N(0, \sigma^2)~~~\mbox{where}~~~
    \sigma^2 \defeq 
    \var_P\left( \frac{A}{\pi(X)}(Y - \mu(X))
    +\mu(X)
    \right).
 \end{equation*}
\end{theorem}
\vspace{-0.7em}
\noindent  Since $\sigma^2$ is the semiparametric efficiency bound for the ATE estimand~\cite{Tsiatis07,Kennedy22}, 
the C-Learner~\eqref{eq:c_learner_def}  achieves the tightest  confident interval  and is optimal in the usual local asymptotic minimax sense~\citep[Theorem 25.21]{van2000asymptotic}.
\vspace{-0.3em}
\paragraph{Double robustness}

By virtue of their first-order correction, standard approaches like one-step estimation and targeting enjoy double robustness: if either of the propensity model or the outcome model is consistent, then the resulting estimator is consistent. We show a similar guarantee for the C-Learner. 
Here, 
we assume that either $\what\pi_{-k,n}$ $\what\mu_{-k,n}^C$ is consistent.
\vspace{-0.9em}
\begin{assumption}[At least one of $\what\pi,\what\mu^C$ is consistent]\label{ass:risk_decay}
For all $k$, the product of the errors for the outcome and propensity
models decays as
\vspace{-0.7em}
    $$\|\what\pi_{-k,n} - \pi\|_{L_2(P)}
    \cdot \|\what\mu^C_{-k,n} - \mu\|_{L_2(P)}=o_P(1).$$
\end{assumption}

Using Assumption~\ref{ass:risk_decay} in place of Assumption~\ref{ass:consistency}, we arrive at the following result; see \cref{sec:double_robustness} for the proof.
\begin{theorem}[C-Learner is doubly robust]
\label{thm:dr}
The C-Learner~\eqref{eq:c_learner_def} is consistent under Assumptions~\ref{ass:overlap}, \ref{ass:bounded_outcomes}, \ref{ass:empirical_proc_mean}, and 
\ref{ass:risk_decay}.
\end{theorem}

\vspace{-0.3em}
\paragraph{``Dual'' C-Learner}
The results in \Cref{thm:asymp} and \Cref{thm:dr} are extended to the ``Dual'' C-Learner as defined in \Cref{eq:c_learner_dual}, and are stated and proved in \Cref{sec:proofs_dual}. 
